% ECONOMICS_PROBLEM: {"id": "D60-517", "title": "Welfare when preferences are inconsistent", "jel_code": "D60", "source_id": "OP-0069"}
% !TeX program = xelatex
\documentclass[11pt,a4paper]{article}
\usepackage[margin=23mm]{geometry}
\usepackage{fontspec}
\setmainfont{Latin Modern Roman}
\usepackage{amsmath,amssymb,mathtools}
\usepackage{xcolor,enumitem,booktabs,longtable,array,xurl,hyperref}
\definecolor{muted}{HTML}{53636C}
\hypersetup{colorlinks=true,urlcolor=blue,linkcolor=blue}
\setlength{\parindent}{0pt}
\setlength{\parskip}{5pt}
\newcommand{\R}{\mathbb R}
\newcommand{\E}{\mathbb E}
\newcommand{\N}{\mathbb N}
\newcommand{\ind}{\mathbf 1}
\newcommand{\Var}{\operatorname{Var}}
\newcommand{\Cov}{\operatorname{Cov}}
\newcommand{\supp}{\operatorname{supp}}
\DeclareMathOperator*{\argmin}{arg\,min}
\DeclareMathOperator*{\argmax}{arg\,max}
\newcommand{\entrymeta}[3]{{\sffamily\small\color{muted}\textbf{\texttt{#1}}\quad|\quad #2\par #3\par}}
\newcommand{\Problemref}[1]{\texttt{#1}}
\newcommand{\JELref}[2]{\texttt{#2} (JEL \texttt{#1})}
\begin{document}
\section*{Welfare when preferences are inconsistent}
\textbf{Problem ID:} \texttt{D60-517}\quad \textbf{JEL:} \texttt{D60}\par
% BEGIN ECONOMICS STATEMENT
\label{problem:OP-0069}
\label{atlas:prob:B9}

\noindent\textbf{Original identifier:} B9 (atlas item 69). \textbf{Classification:} Normative identification and robust policy research program. \textbf{Original tractability label:} Unknown (editorial, not a complexity result).

\subsection*{Definitions and assumptions}

Let $X$ be a finite set of economically relevant allocations and let $C(B,f,t)$ be an observed choice correspondence from feasible menu $B\subseteq X$ under frame $f$ and date $t$. A behavioral data set records available alternatives as well as choices. Supply a nonempty family $\mathcal W$ of finite real-valued welfare functionals with declared cardinal and interpersonal normalizations, and justify which decision conditions are welfare-relevant. Each $W_w(\pi;P)$ must be finite on the declared nonempty model class $\mathcal P$; purely ordinal relations require a separate comparison definition. For policies $\pi\in\Pi$, let $P_\pi$ be the identified or bounded allocation distribution, including implementation costs and effects on others. Choice inconsistency alone neither identifies $\mathcal W$ nor licenses an arbitrary paternalistic criterion.

\subsection*{Formal research question}

Characterize the policy ranking robust to both positive-model and welfare ambiguity:
\[
\pi\succeq_R\pi'\quad\Longleftrightarrow\quad
\inf_{w\in\mathcal W,\;P\in\mathcal P}\{W_w(\pi;P)-W_w(\pi';P)\}\geq0,
\]
where $\mathcal P$ is a declared set of models consistent with the data and maintained restrictions. Determine conditions for nontrivial robust comparisons, and characterize the pairs for which rankings reverse across admissible evaluators. For specified policies, identify additional evidence that reduces $\mathcal P$ or narrows a substantively justified $\mathcal W$ enough to determine a ranking.

\subsection*{Interpretation and scope}

This problem contains a normative input that statistics cannot uniquely recover without assumptions. Potential evaluators include informed deliberative choice, an identified long-run self, experienced utility, or an explicitly justified choice-based relation excluding certain frames. Each has different informational and ethical commitments; the exercise must state them rather than silently choosing the frame whose recommendation is preferred. Even a complete description of all observed choices may leave welfare incomplete. Policy effects can be consistently estimated while welfare comparisons remain indeterminate. The robust relation is a preorder. It induces a partial order only on equivalence classes under mutual robust comparison; incomparability remains an informative outcome. Numerical aggregators must define interpersonal weights and normalization when policy affects multiple people. Report which recommendations are stable and which reverse, separating causal uncertainty from normative disagreement. A proposed experiment can test factual assumptions about attention or comprehension, but cannot by itself prove that one normative evaluator is uniquely correct.

\subsection*{Source audit and status}

[S1] provides a behavioral choice-based welfare framework. The ambiguous Kőszegi--Rabin (2007) citation is resolved to the welfare paper [S2]; their distinct 2007 risk-attitudes article would have a different role. [S3] supplies a pragmatic policy perspective. All selected works are verified. No source establishes a uniquely correct welfare preference for every inconsistent chooser. The robust comparison target is an editorial synthesis, and its normative ambiguity should not be mislabeled solely as an unresolved mathematical theorem.

\subsection*{References}

\begin{enumerate}[label={[S\arabic*]},leftmargin=*]

\item B. Douglas Bernheim and Antonio Rangel (2009). \emph{Beyond Revealed Preference: Choice-Theoretic Foundations for Behavioral Welfare Economics}. Quarterly Journal of Economics 124(1), 51--104. \url{https://doi.org/10.1162/qjec.2009.124.1.51}\par \textit{Source role:} Choice-based welfare framework. \textit{Locator:} Abstract and article metadata.

\item Botond Kőszegi and Matthew Rabin (2007). \emph{Mistakes in Choice-Based Welfare Analysis}. American Economic Review 97(2), 477--481. \url{https://www.aeaweb.org/articles?id=10.1257/aer.97.2.477}\par \textit{Source role:} Welfare critique; selected interpretation of ambiguous author-year. \textit{Locator:} Abstract and article metadata.

\item Raj Chetty (2015). \emph{Behavioral Economics and Public Policy: A Pragmatic Perspective}. American Economic Review 105(5), 1--33. \url{https://www.aeaweb.org/articles?id=10.1257/aer.p20151108}\par \textit{Source role:} Policy and sufficient-statistic perspective. \textit{Locator:} Abstract and article metadata.

\end{enumerate}

\subsection*{Common conventions: Source mathematical conventions}

Unless an entry states otherwise, agent and firm sets are finite. State and action spaces are measurable, strategies and outcome maps are measurable, probability laws are specified, and expectations exist for the targets under discussion. Infinite-horizon discount factors lie in $(0,1)$ and discounted objectives are finite. Norms, tolerances and complexity input sizes must be specified for computational comparisons. Constants or primitives that appear locally are inputs, not empirical facts asserted by this catalogue.

\textbf{Identification.} A statistical model $m\in\mathcal M$ implies an observable law $P_m$ and target $\theta(m)$. At law $P$, its identified set is
\[
\Theta(P;\mathcal M)=\{\theta(m):m\in\mathcal M,\ P_m=P\}.
\]
Point identification holds when this set is a singleton, or equivalently when $P_{m_1}=P_{m_2}$ implies $\theta(m_1)=\theta(m_2)$ for all compatible models. Sharp bounds mean bounds attained, or approached when the boundary is not attained, by compatible models. Writing this set is a definition, not a solution: the substantive task is to establish informative restrictions and derive or estimate the set. An unrestricted-confounding model may give only trivial bounds.

\textbf{Inference.} Identification, estimability and valid uncertainty quantification are separate questions. Fix $\alpha\in(0,1)$. A confidence procedure $C_n$ over a declared law class $\mathcal P$ has asymptotic uniform coverage at level $1-\alpha$ when
\[
\liminf_{n\to\infty}\inf_{P\in\mathcal P}\Pr_P\{\theta(P)\in C_n\}\geq1-\alpha.
\]
For set-identified targets the coverage event must be defined separately, for example coverage of the full identified set. If an entry uses identification-indexed models instead of uniquely indexed laws, the same coverage convention is applied over those models. Pointwise limits do not establish uniform validity, and asymptotic validity does not guarantee finite-sample precision. Sample sizes, dependence structures and weak-identification sequences are part of each application.

\textbf{Counterfactuals.} $Y_i(a)$ is a potential outcome under a specified intervention, including its timing, implementation and versions. It is not $Y_i$ conditional on voluntarily choosing $a$. A full intervention vector permits interference; shorthand scalar treatment notation does not silently impose no interference. Assignment, selection, attrition, anticipation and measurement must be specified. A claim about an instrument's local effect is not automatically a population policy effect.

\textbf{Economic equilibrium.} A counterfactual model includes best responses, constraints, feasibility, market clearing, state transitions and expectations consistent with the stated information. When several equilibria exist, either a declared selector is an input or the target is set-valued. Comparisons across policy regimes must use the stated selection convention. Reduced-form relationships are not presumed invariant after an equilibrium-changing intervention.

\textbf{Optimization and welfare.} Optimization is over the stated feasible set. If attainment is not established, $\sup$ or $\inf$ and $\varepsilon$-optimal choices replace an asserted maximizer or minimizer. For example, with value $V=\sup_{a\in\mathcal A}W(a)<\infty$, an $\varepsilon$-optimal set is $\{a\in\mathcal A:W(a)\geq V-\varepsilon\}$ for $\varepsilon>0$. Benefits, costs, risk and health or environmental outcomes can be aggregated only using declared units and conversion weights. Interpersonal utility weights, the treatment of fiscal transfers and foreign profits, discounting, and the behavioral welfare criterion are explicit inputs. A comparative-static derivative additionally requires existence and appropriate differentiability of the selected counterfactual.

\textbf{Sources.} Local labels [S1], [S2], and so forth restart in every question. Locators identify the relevant abstract, model, section or result at the level actually checked; a bibliographic or abstract-level verification is not represented as inspection of an entire proof. Working papers are distinguished from later publications and changing versions. Source summaries are paraphrased. All URL checks and editorial formulations refer to the audit date stated above.
% END ECONOMICS STATEMENT
\end{document}
